\section{Considerações Finais}
\label{sec:consideracoes}

Este relatório documentou três experimentos de adaptação do \texttt{Qwen2.5-0.5B} ao
domínio público piauiense. A Questão 01, concluída integralmente, confirma sua
hipótese de forma robusta: o pré-treino continuado sobre o DOMPI-2025 reduziu a
perplexidade de 21,78 para 2,09 (queda de 90,4\%) e elevou a acurácia de previsão de
tokens de 42,25\% para 82,00\%, sem sinais de \textit{overfitting} e com clara
transformação qualitativa das gerações, que passaram de texto genérico para linguagem
administrativa autêntica do Piauí.

As Questões 02 (SFT completo) e 03 (LoRA/QLoRA) tiveram a metodologia, a geração dos
3.111 pares de instrução/resposta e a avaliação \textit{baseline} concluídas com
números reais, mas o treinamento foi interrompido por restrição de tempo de hardware
na tentativa local. Optou-se por reportar esse estado de forma transparente, sem
estimar métricas pós-treino. Os dados já evidenciam o principal trade-off entre as
duas técnicas: o LoRA/QLoRA atualiza apenas 0,44\% dos parâmetros e estimou treino
cerca de 20 vezes mais rápido que o ajuste fino completo.

\subsection{Caminho para conclusão das Questões 02 e 03}

Para fechar as duas questões pendentes, basta reexecutar os notebooks na GPU T4 do
Kaggle (o mesmo ambiente que viabilizou a Questão 01 em $\approx$4h40), aplicando os
ajustes já identificados:

\begin{enumerate}[leftmargin=2.2em]
  \item aplicar as três melhorias de qualidade dos pares de QA da
        Seção~\ref{sec:q2} (deduplicação, filtro de idioma, máscara de \textit{loss})
        e as perguntas de avaliação com nomes reais de professores;
  \item aplicar a correção do modo LoRA puro da Seção~\ref{sec:q3}
        (\texttt{enable\_input\_require\_grads});
  \item executar o treino, coletar as métricas pós-treino e preencher as células
        ``depois'' das Tabelas~\ref{tab:q2-baseline}, \ref{tab:q3-baseline}
        e~\ref{tab:comparativo}.
\end{enumerate}

Como o conjunto docentesDC é muito menor que o DOMPI-2025, estima-se que cada treino
leve apenas dezenas de minutos na T4, dominado pelo custo fixo de carregamento e
avaliação. A estrutura deste relatório já está preparada para receber esses números.

\subsection{Trabalhos futuros}

O trabalho da disciplina prevê ainda três frentes não cobertas por este relatório,
que dão sequência natural ao que foi construído: \textbf{(4)} destilação de
conhecimento de um modelo professor para um aluno; \textbf{(5)} uma aplicação de RAG
(\textit{Retrieval-Augmented Generation}) sobre o DOMPI-2025 ou o docentesDC, capaz
de recuperar fatos pontuais (CNPJ, CEP, e-mail) que o pré-treino não memoriza; e
\textbf{(6)} a inclusão de camadas de \textit{guardrails} sobre um dos modelos
desenvolvidos. A frente de RAG, em particular, complementa diretamente a limitação de
memorização factual observada no \textit{benchmark} da Questão 01.
